% T23 -- Capteurs et conditionnement du signal -- banque d'exercices
%
% COMPLEMENTARITE AVEC LA FICHE fiches/T23_Capteurs_Conditionnement.tex
% Les trois Questions Types de la fiche sont :
%   QT1 = identification d'un capteur PT100 Jumo + coefficient c = 3/95 = 0,0316 V/degC
%         (attribuee a "E2 BAC3 2026") ;
%   QT2 = pont de Wheatstone 4 jauges, R0 = 120 Ohm, Vcc = 1 V, dR = 0,0153 Ohm,
%         Vout = 31,9 uV et gain Ka = 1,56e5. C'est EXACTEMENT le couple de questions
%         58/59 du sujet 2025 (epreuve ecrite disciplinaire appliquee) ;
%   QT3 = capteur inductif f = N.Z.n/60 = 96 Hz sur Z = 24 dents a 240 tr/min. C'est
%         EXACTEMENT la question 15 du sujet E1 2020.
% CES TROIS TRAVAUX INTELLECTUELS NE SONT PAS REPRIS ICI :
%   - aucun exercice ne recalcule Vout = (Vcc/4).(dR/R0) sur les donnees 120 Ohm / 1 V /
%     0,0153 Ohm de 2025, ni le gain 1,56e5 associe ;
%   - aucun exercice ne calcule f = N.Z.n/60 pour un capteur inductif ;
%   - aucun exercice ne calcule un simple coefficient de pente (Vmax-Vmin)/(Tmax-Tmin).
% Le pont de Wheatstone EST traite (ex. 5), mais par un angle que la fiche n'aborde pas :
% la VALIDITE de l'hypothese dR << R0 sur une PT100 (ou dR/R0 = 0,58 : la linearisation
% est fausse de 29 %), la ou la fiche applique la formule linearisee sans jamais la
% discuter. Aucun chiffre de 2025 n'y est reemploye.
%
% SOURCES DE CHAQUE EXERCICE (toutes relues dans les .md de conversion Marker sous
% docs/CAPET_Markdown-20260629T134311Z-3-001/CAPET_Markdown/docs/CAPET_SII_IE/06_PDFs_Originaux/)
%
%   - Ex. 1 (docTech) : 2024, "EPREUVE ECRITE DISCIPLINAIRE", option Ingenierie
%     electrique (code EDE ENE 1 ; nom lu en tete de E2_Disciplinaire_IE_2024.md l.1-13),
%     question D1. Le tableau du docTech est la "Figure 13 : caracteristiques en fonction
%     du type de thermocouple", TEXTUELLEMENT presente dans
%     E2_Disciplinaire_IE_2024.md l.463-475 (tableau markdown l.466-472) : les cinq lignes
%     T / K / N / J / E, leurs plages, alliages, tolerances, prix et remarques de fiabilite
%     sont recopies sans rien y ajouter. Contexte (thermocouple, principe Seebeck) : l.445-461.
%     Plage de mesure -20 a 200 degC : l.687 et DT10 l.937-950.
%     Corrige officiel de D1 : Rapport_Jury_2024.md l.313.
%     ATTENTION : le rapport ecrit "une tolerance pour des temperatures autour de 40 degC de
%     40x2,5/100 = 1 degC" ; or la figure 13 donne au type J une tolerance de +/- 1,5 (et non
%     2,5, qui est celle des types K et N). Le "2,5" du rapport ne correspond pas au type J
%     qu'il conclut. Le corrige ci-dessous reproduit donc la demarche du jury (comparer la
%     tolerance a l'exigence de l'application) SANS reprendre ce calcul incoherent.
%
%   - Ex. 2 : 2023, "EPREUVE ECRITE DISCIPLINAIRE APPLIQUEE" (nom lu en tete de
%     E2_Appliquee_Commune_2023.md, ligne 13 ; epreuve commune aux cinq options),
%     questions 46, 47 et 50. Enonce : E2_Appliquee_Commune_2023.md l.607-629.
%     Debattement mecanique -23 deg a +23 deg (l.607), signal 4-20 mA issu d'une sonde a
%     effet Hall (l.611), trajet de mesurage SP1 = ES11 et SP2 = ES12 (l.623-627).
%     Corrige officiel : Rapport_Jury_2023.md l.912-918 (Q46 : "Une rupture de ligne ou une
%     panne de courant donnent I voisin de zero qui devient alors un indice de surveillance.")
%     La question 3 (reprogrammation SP2' = 0 deg) s'appuie sur la question 50 du meme sujet
%     (enonce l.655-657 ; plage utile -23 deg a 0 deg pour COUV <= 50 %) et sur le corrige
%     Rapport_Jury_2023.md l.931.
%
%   - Ex. 3 : 2024, meme epreuve que l'ex. 1, questions E1, E2 et E3.
%     Enonce : E2_Disciplinaire_IE_2024.md l.700-755. Valeurs R1 = 1 kOhm et R2 = 834,4 Ohm :
%     l.738 (texte de la figure 22). Chaine des trois gains G1 / G2 = -100 / G3 = -1 : l.724
%     (texte de la figure 21). Relation V1 = VR2.(2R1+R2)/R2 : l.752.
%     Corrige officiel : Rapport_Jury_2024.md l.435-441 (E1 : eAB = VR2 ; E2 : G1 = 3,397 ;
%     E3 : V3 = 340.eAB). Recalcul independant : (2000+834,4)/834,4 = 3,39693 et
%     3,39693 x 100 = 339,69, arrondi a 340 par le jury -- confirme.
%
%   - Ex. 4 : 2024, meme epreuve, questions E5 et E10 (l'etude detaillee de la boucle de
%     conversion, questions E6 a E9, n'est pas reprise : elle exige la figure 24 qui est une
%     image omise par la conversion). Donnees : Ve = 1,34 V et V4 = 337,7.eAB + 1,34 (l.763) ;
%     R5 = 250 Ohm (l.797) ; eABmini = -1 mV et eABmaxi = 10,8 mV pour -20 degC et 200 degC
%     (DT10, tableau l.945-950 et l.1141) ; sortie Is 4-20 mA (l.687).
%     Corrige officiel : Rapport_Jury_2024.md l.451-453 (V4 = 1 V et 5 V) et l.467
%     (Is = 4 mA et 20 mA, "On obtient bien le profil de courant que doit delivrer le
%     transmetteur soit 4 mA pour -20 degC et 20 mA pour 200 degC").
%     Recalcul independant : 337,7 x (-0,001) + 1,34 = 1,0023 V et
%     337,7 x 0,0108 + 1,34 = 4,98716 V, arrondis a 1 V et 5 V par le jury -- l'ecart vient
%     de l'arrondi du coefficient 337,7 (la valeur exacte est 4/0,0118 = 338,98). Signale.
%
%   - Ex. 5 : EXERCICE INEDIT (aucune reference de sujet). Construit sur la formule de pont
%     de Wheatstone quart de pont et sur les caracteristiques normalisees de la PT100
%     (R0 = 100 Ohm a 0 degC, alpha = 3,85e-3 /degC : valeurs de la norme CEI 60751, deja
%     presentes dans la fiche T23 bloc 1.2, non controversees, non re-verifiees).
%     Toutes les AN sont recalculees. Interet pedagogique = MISE EN DEFAUT de l'hypothese
%     dR << R0 : a 150 degC, dR/R0 = 0,5775 et la formule linearisee surestime de 28,9 %.
%     Verifications : dR(150) = 100 x 3,85e-3 x 150 = 57,75 Ohm ;
%       Vout_lin  = (5/4) x 0,5775            = 0,721875 V ;
%       Vout_exact= (5/2) x 57,75/257,75      = 0,560136 V ;
%       erreur    = 0,721875/0,560136 - 1     = +28,875 % = 1 + 57,75/200 - 1. OK.
%       K = 3/0,560136 = 5,3558 ; K_faux = 3/0,721875 = 4,1558 (-22,4 %).
%
%   - Ex. 6 : 2022, "EPREUVE ECRITE DISCIPLINAIRE", option Ingenierie electrique
%     (code EDE ENE 1 ; nom lu en tete de E2_Disciplinaire_IE_2022.md l.1-5), question 24.
%     Enonce : E2_Disciplinaire_IE_2022.md l.445-467. Tension de reference du CAN 3,5 V
%     (l.455) ; CAN 12 bits du PIC 18F8723 (l.457 et Q23) ; pont diviseur R4/R5 de rapport 2
%     entre le module Peltier et l'entree ADC2 ; alarme sur 0xAF = 175 (l.469).
%     Corrige officiel : Rapport_Jury_2022.md l.305-319 (quantum = 0,85 mV/LSB ;
%     N = 1170 pour UPELTIER = 2 V ; ADC2 = 0,15 V et UPELTIER = 0,3 V pour 175).
%     Recalcul independant : 3,5/4095 = 0,854701 mV/LSB ; 1/(3,5/4095) = 1170,0 exactement ;
%     175 x 3,5/4095 = 0,149573 V, x2 = 0,299145 V. Confirme.
%     ATTENTION : le corrige du jury conclut "donc inferieure au 0,4 V minimum impose sur le
%     diagramme des exigences" alors que l'enonce du sujet (l.523) ecrit "Cette valeur de la
%     tension ne doit pas etre inferieure a 0,9 V". Les deux seuils divergent (le diagramme
%     des exigences est une image omise par la conversion, impossible de trancher). Le
%     corrige ci-dessous utilise le seuil de 0,9 V de l'ENONCE et signale la divergence : la
%     conclusion (alarme justifiee) est la meme dans les deux cas puisque 0,3 < 0,4 < 0,9.
%
% CITATIONS DE RAPPORT DE JURY -- verbatim exacts verifies par grep, avec ligne :
%   - Ex. 1  : Rapport_Jury_2024.md l.515 (analyse globale, epreuve ecrite disciplinaire).
%   - Ex. 5  : Rapport_Jury_2025.md l.1024 (commentaire de la partie 5).
%   - Ex. 6  : Rapport_Jury_2022.md l.389 (commentaires de l'epreuve ecrite disciplinaire).
% Aucune autre citation n'est produite. Aucune valeur n'est attribuee a un sujet sans avoir
% ete relue ligne a ligne dans le .md correspondant.
\expandafter\ifx\csname CAPETINCLUDEMODE\endcsname\relax
\documentclass[11pt,a4paper]{article}
% preamble.tex — packages et commandes communes a toutes les fiches CAPET SII IE
\input{glyphtounicode}
\pdfgentounicode=1
\usepackage{cmap}
\usepackage[T1]{fontenc}
\usepackage[utf8]{inputenc}
\usepackage[french]{babel}
% babel[french] active a \begin{document} des raccourcis de ponctuation (:;!?)
% qui cassent la syntaxe to[...] de circuitikz (erreur "+ or - expected") ;
% on les desactive, ce qui ne change que l'espacement fin avant ces signes.
\AtBeginDocument{\shorthandoff{:;!?}}
\usepackage[a4paper, top=1.6cm, bottom=1.6cm, left=1.4cm, right=1.4cm]{geometry}
\setlength{\headheight}{14.5pt}
\usepackage{amsmath, amssymb, amsfonts}
\usepackage{siunitx}
\sisetup{locale=FR, output-decimal-marker={,}, inter-unit-product=\ensuremath{{}\cdot{}},
         group-minimum-digits=4}
% Unites non fournies par siunitx, declarees ici une fois pour toutes (fiches + exercices).
% Les gardes \ifdefined evitent tout conflit avec une redeclaration locale restee dans un
% fichier : ne pas retirer les gardes, ne pas redeclarer ces unites dans les fichiers.
\DeclareSIUnit{\jour}{j}
\ifdefined\var\else\DeclareSIUnit{\var}{var}\fi   % puissance reactive
\ifdefined\vad\else\DeclareSIUnit{\vad}{VAD}\fi   % puissance deformante
\ifdefined\tr\else\DeclareSIUnit{\tr}{tr}\fi      % tour (tr/min)
\usepackage{array, booktabs, multirow, tabularx}
\usepackage{xcolor}
\usepackage[most, breakable]{tcolorbox}
\usepackage{multicol}
\usepackage{enumitem}
\usepackage{mathtools}
\usepackage{lmodern}
% Refonte visuelle 2026 (direction "moderne coloré") : Lato pour le texte,
% maths sans-serif assorties via newtxsf (charge apres amsmath/amssymb).
\usepackage[default]{lato}
\usepackage{newtxsf}
\usepackage{microtype}
\usepackage{fancyhdr}
\usepackage{tikz}
\usetikzlibrary{positioning,arrows.meta}
\usepackage{circuitikz}
\usepackage{lastpage}

% ── Couleurs par domaine (palette modernisée 2026) ────────────────────
\definecolor{colElec}{RGB}{29, 78, 216}        % Électrotechnique — bleu vif
\definecolor{colElecLight}{RGB}{219, 234, 254} % fond bleu clair
\definecolor{colAuto}{RGB}{21, 128, 61}        % Automatique — vert
\definecolor{colMeca}{RGB}{194, 65, 12}        % Mécanique — orange
\definecolor{colTherm}{RGB}{185, 28, 28}       % Thermique — rouge
\definecolor{colQT}{RGB}{232, 244, 233}        % (hérité, plus utilisé par questiontype)
\definecolor{colJury}{RGB}{255, 248, 225}      % fond jury (crème)
\definecolor{colJuryFrame}{RGB}{202, 138, 4}   % bordure/bandeau jury (ambre)
\definecolor{colGrayLight}{RGB}{245, 245, 247} % fond gris neutre
\definecolor{colSeparator}{RGB}{180, 180, 180} % gris séparateur
\definecolor{colExo}{RGB}{8, 145, 178}          % Banque d'exercices — cyan (distinct des fiches)
% ── Encre et gris de la refonte ───────────────────────────────────────
\definecolor{colInk}{RGB}{23, 32, 42}          % texte courant (quasi-noir)
\definecolor{colSlate}{RGB}{83, 98, 118}       % gris ardoise (méta, en-têtes)
\definecolor{colRule}{RGB}{203, 213, 225}      % filets clairs
\definecolor{colJuryBg}{RGB}{255, 248, 225}    % fond crème jury
\definecolor{colPresence}{RGB}{83, 98, 118}    % bloc 5 (présence dans les sujets)
\AtBeginDocument{\color{colInk}}

% ── Style encadré exercice + corrigé (banque d'exercices) ─────────────
\tcbset{
  exoBox/.style={
    breakable,
    enhanced,
    colback=colExo!8!white,
    colframe=colExo,
    colbacktitle=colExo,
    coltitle=white,
    fonttitle=\bfseries\small,
    left=4pt, right=4pt, top=3pt, bottom=3pt,
    before skip=5pt, after skip=5pt,
  }
}

% ── Bascule enonce / corrige (banque d'exercices) ─────────────────────
% \CAPETCORRIGE est defini en ligne de commande (pdflatex "\def\CAPETCORRIGE{1}\input{...}")
% pour compiler la variante corriges d'un meme fichier source ; par defaut (non defini),
% la compilation produit la variante enonces. \CAPETSHOW{enonce}{corrige} choisit la branche.
%
% TROIS MODES (refonte 2026-09-20, manuel de revision) :
%   - aucun drapeau   -> imprime l'enonce            (variante enonces, INCHANGE)
%   - \CAPETCORRIGE   -> imprime le corrige          (variante corriges, INCHANGE)
%   - \CAPETMANUEL    -> imprime l'enonce ET met le corrige de cote pour l'annexe
%
% Le mode manuel stocke chaque corrige dans une macro (\csgdef), sans ecriture
% fichier ni re-tokenisation : le contenu (circuitikz, accents, %) traverse intact.
% L'annexe les rejoue ensuite dans l'ordre via \CAPETANNEXE.
\usepackage{etoolbox}

\newcounter{CAPETexo}          % numero continu d'exercice dans le manuel
\setcounter{CAPETexo}{0}
\newcounter{CAPETannexe}       % compteur de relecture en annexe

% Titre courant de la notion en cours (pose par \CAPETNOTION), sert d'intertitre annexe
\newcommand{\CAPETnotioncode}{}
\newcommand{\CAPETnotiontitre}{}

% \CAPETNOTION{code}{titre} : ouvre une notion dans le manuel (memorise le libelle,
% et marque le rang du premier exercice de la notion pour l'intertitre d'annexe).
\newcommand{\CAPETNOTION}[2]{%
  \renewcommand{\CAPETnotioncode}{#1}%
  \renewcommand{\CAPETnotiontitre}{#2}%
  \csgdef{CAPETnotionstart\the\numexpr\value{CAPETexo}+1\relax}{#1 -- #2}%
}

% ── \degree robuste ───────────────────────────────────────────────────
% Certains fichiers d'exercices utilisent \degree. Ni le preambule ni siunitx
% ne le definissent, et \providecommand local ne suffit pas : sous hyperref,
% \degree existe deja comme commande factice (\ERROR), donc \providecommand
% s'abstient et l'appel echoue a la compilation du manuel (erreur invisible
% en compilation isolee). On le (re)definit donc inconditionnellement ici,
% apres le chargement des packages, pour tout le corpus.
\AtBeginDocument{\providecommand{\degree}{}\renewcommand{\degree}{\ensuremath{^\circ}}}

% Libelle decoratif : les fichiers d'exercices utilisent aussi \CAPETSHOW dans
% leur \fancyhead et leur bloc de titre, pour afficher "(enonces)" ou "(corriges)".
% Ces appels-la ne sont PAS des exercices : en mode manuel ils ne doivent ni
% consommer un numero, ni alimenter l'annexe. On les reconnait a leur premier
% argument, strictement egal au libelle ci-dessous.
\newcommand{\CAPETLIBELLEENONCES}{(\'enonc\'es)}

\ifdefined\CAPETMANUEL
  % --- Mode manuel : enonce en place, corrige differe vers l'annexe ---
  \newcommand{\CAPETSHOW}[2]{%
    \def\CAPETargun{#1}%
    \ifx\CAPETargun\CAPETLIBELLEENONCES
      #1%            appel decoratif : rendu tel quel, hors numerotation
    \else
      \CAPETSHOWEXO{#1}{#2}%
    \fi
  }
  \newcommand{\CAPETSHOWEXO}[2]{%
    \stepcounter{CAPETexo}%
    \csgdef{CAPETcorrige\theCAPETexo}{#2}%
    \edef\CAPETtmpnum{\theCAPETexo}%
    #1%
    \par\vspace{2pt}%
    {\footnotesize\itshape\color{colExo}%
     \hyperref[capet:cor:\CAPETtmpnum]{Corrige de l'exercice \CAPETtmpnum\ $\rightarrow$ p.~\pageref{capet:cor:\CAPETtmpnum}}%
     \label{capet:eno:\CAPETtmpnum}}%
  }
  % \CAPETANNEXE : rejoue tous les corriges captures, dans l'ordre, avec intertitres
  \newcommand{\CAPETANNEXE}{%
    \setcounter{CAPETannexe}{0}%
    \loop\ifnum\value{CAPETannexe}<\value{CAPETexo}%
      \stepcounter{CAPETannexe}%
      \edef\CAPETcurnum{\theCAPETannexe}%
      \ifcsname CAPETnotionstart\CAPETcurnum\endcsname
        \par\medskip
        {\normalsize\bfseries\color{colExo}\csuse{CAPETnotionstart\CAPETcurnum}}%
        \par\nobreak\smallskip
      \fi
      \begin{tcolorbox}[exoBox, breakable,
        title={Corrige de l'exercice \CAPETcurnum\
               \hfill\normalfont\footnotesize
               \hyperref[capet:eno:\CAPETcurnum]{enonce p.~\pageref{capet:eno:\CAPETcurnum}}}]%
      \label{capet:cor:\CAPETcurnum}%
      \csuse{CAPETcorrige\CAPETcurnum}%
      \end{tcolorbox}%
    \repeat
  }
\else
  % --- Modes historiques, comportement strictement inchange ---
  \newcommand{\CAPETSHOW}[2]{\ifdefined\CAPETCORRIGE #2\else #1\fi}
  \newcommand{\CAPETANNEXE}{}
\fi

% ── Difficulté 1 à 3 (étoiles pleines, amssymb) ───────────────────────
\newcommand{\difficulte}[1]{%
  \ifcase#1\relax
  \or $\bigstar$\;niveau 1/3%
  \or $\bigstar\bigstar$\;niveau 2/3%
  \or $\bigstar\bigstar\bigstar$\;niveau 3/3%
  \fi
}

% ── Style encadré extrait de document technique (banque d'exercices) ──
\tcbset{
  docTech/.style={
    breakable,
    enhanced,
    colback=colGrayLight,
    colframe=colSeparator,
    fonttitle=\bfseries\small\itshape,
    coltitle=black,
    colbacktitle=colGrayLight,
    left=4pt, right=4pt, top=3pt, bottom=3pt,
    before skip=4pt, after skip=4pt,
  }
}

% ── Style tcolorbox de section principale ────────────────────────────
% Refonte 2026 : bandeau titre coloré arrondi + fond très pâle, sans cadre.
\tcbset{
  ficheSection/.style={
    breakable, enhanced jigsaw,
    frame hidden, boxrule=0pt, arc=1.8mm,
    colback=#1!4!white,
    colbacktitle=#1,
    coltitle=white,
    fonttitle=\sffamily\bfseries,
    left=7pt, right=7pt, top=5pt, bottom=5pt,
    toptitle=3pt, bottomtitle=3pt,
    before skip=10pt, after skip=8pt,
  }
}

% ── Style encadré question type ──────────────────────────────────────
% Carte blanche, bandeau sombre, liseré très clair.
\tcbset{
  questiontype/.style={
    breakable, enhanced jigsaw,
    colback=white,
    frame hidden, boxrule=0pt, arc=1.8mm,
    colbacktitle=colElec!78!black,
    coltitle=white,
    fonttitle=\sffamily\bfseries\small,
    left=7pt, right=7pt, top=4pt, bottom=4pt,
    toptitle=2.5pt, bottomtitle=2.5pt,
    before skip=8pt, after skip=8pt,
    borderline={0.6pt}{0pt}{colElec!12!white},
  }
}

% ── Style encadré jury ───────────────────────────────────────────────
% Bandeau ambre, fond crème, sans cadre.
\tcbset{
  jury/.style={
    breakable, enhanced jigsaw,
    colback=colJuryBg,
    frame hidden, boxrule=0pt, arc=1.8mm,
    colbacktitle=colJuryFrame,
    coltitle=white,
    fonttitle=\sffamily\bfseries,
    left=7pt, right=7pt, top=4pt, bottom=4pt,
    toptitle=3pt, bottomtitle=3pt,
    before skip=10pt, after skip=8pt,
  }
}

% ── En-tête / pied de page communs (refonte 2026) ─────────────────────
% #1 code fiche, #2 titre court, #3 couleur domaine
\newcommand{\ficheHeaders}[3]{%
  \fancyhead[L]{\small\sffamily{\bfseries\textcolor{#3}{#1}}\hspace{0.7em}\textcolor{colSlate}{#2}}%
  \fancyhead[R]{\small\sffamily\textcolor{colSlate}{CAPET SII IE\, --- \,Session 2026}}%
  \fancyfoot[C]{\small\sffamily\textcolor{colSlate}{\thepage\,/\,\pageref{LastPage}}}%
  \fancyhead[C]{}\fancyfoot[L]{}\fancyfoot[R]{}%
}

% ── Bloc en-tête de fiche (BLOC 0, refonte 2026) ──────────────────────
% #1 code  #2 titre  #3 couleur domaine  #4 domaine  #5 ligne présence  #6 sous-thèmes
\newcommand{\ficheEntete}[6]{%
  \begin{tcolorbox}[enhanced, breakable,
    frame hidden, boxrule=0pt, arc=2.5mm,
    interior style={left color=#3, right color=#3!75!black},
    left=9pt, right=9pt, top=7pt, bottom=7pt,
    before skip=0pt, after skip=6pt]
    {\color{white}\sffamily\bfseries\LARGE #1\quad #2}\\[4pt]
    {\color{#3!22!white}\sffamily\small \MakeUppercase{#4}\quad\textbar\quad #5}\\[1pt]
    {\color{#3!35!white}\sffamily\footnotesize Sous-thèmes : #6}
  \end{tcolorbox}%
}

% ── Mise en page ─────────────────────────────────────────────────────
\pagestyle{fancy}
\fancyhf{}
\renewcommand{\headrulewidth}{0.5pt}
\renewcommand{\footrulewidth}{0pt}
\AtBeginDocument{%
  \renewcommand{\headrule}{{\color{colRule}\hrule width\headwidth height 0.5pt}}%
}

\setlist[itemize]{leftmargin=1.2em, itemsep=1pt, topsep=2pt, parsep=0pt}
\setlist[enumerate]{leftmargin=1.2em, itemsep=3pt, topsep=2pt, parsep=0pt}

\setlength{\parskip}{2pt}
\setlength{\columnsep}{10pt}

\def\CAPETBEGINFILE{\begin{document}}
\def\CAPETENDFILE{\end{document}}
\else
\def\CAPETBEGINFILE{}
\def\CAPETENDFILE{}
\fi

% \degree n'est pas defini par le preambule partage (ni par gensymb, non charge) :
% on le fournit localement, sans ecraser une eventuelle definition future.
\providecommand{\degree}{\ensuremath{^\circ}}

\fancyhead[L]{\textbf{\textcolor{colExo}{T23}} - Exercices - Capteurs et conditionnement}
\fancyhead[R]{\small\textit{Banque d'exercices \CAPETSHOW{(\'enonc\'es)}{(corrig\'es)}}}
\fancyfoot[C]{\small\thepage\,/\,\pageref{LastPage}}

\CAPETBEGINFILE

\begin{tcolorbox}[
  enhanced, breakable,
  colback=colExo,
  colframe=colExo,
  left=6pt, right=6pt, top=5pt, bottom=5pt,
  before skip=0pt, after skip=4pt
]
  {\color{white}\bfseries\large T23 \quad Capteurs et conditionnement du signal}\\[2pt]
  {\color{white}\small Banque d'exercices -- 6 exercices \CAPETSHOW{(\'enonc\'es)}{(corrig\'es)}}
\end{tcolorbox}

% =====================================================================
% Exercice 1 -- docTech -- 2024, question D1
% =====================================================================
\begin{tcolorbox}[exoBox,
  title={Exercice 1 -- \difficulte{1} \hfill\normalfont\footnotesize 2024, \'epreuve \'ecrite disciplinaire (option IE), question D1 (figure 13)}
]
\CAPETSHOW{
La temp\'erature $T_{ARM}$ \`a l'int\'erieur de l'armoire \'electrique d'un extracteur d'air industriel doit \^etre surveill\'ee afin de piloter un ventilateur de refroidissement. Le capteur retenu est un \textbf{thermocouple}.

\emph{Principe de fonctionnement rappel\'e par le sujet} : si l'on r\'eunit \`a une extr\'emit\'e (jonction) deux fils m\'etalliques de natures diff\'erentes et que l'on \'el\`eve la temp\'erature de cette extr\'emit\'e, il appara\^it une tension $e_{AB}$ aux extr\'emit\'es rest\'ees libres.

\begin{center}
\begin{tikzpicture}[>=Stealth, scale=0.9, every node/.style={scale=0.9}]
  \draw[thick] (4.0,0.0) -- (1.2,0.55);
  \draw[thick] (4.0,0.0) -- (1.2,-0.55);
  \fill (4.0,0.0) circle (2.2pt);
  \node[right, font=\scriptsize, align=left] at (4.15,0.0) {soudure\\\`a la jonction\\($T_{ARM}$)};
  \node[above, font=\scriptsize] at (2.6,0.36) {M\'etal A};
  \node[below, font=\scriptsize] at (2.6,-0.36) {M\'etal B};
  \draw[thick] (1.2,0.55) -- (0.55,0.55);
  \draw[thick] (1.2,-0.55) -- (0.55,-0.55);
  \fill (0.55,0.55) circle (1.6pt);
  \fill (0.55,-0.55) circle (1.6pt);
  \draw[<->, thin] (0.28,0.55) -- (0.28,-0.55);
  \node[left, font=\scriptsize] at (0.22,0.0) {$e_{AB}$};
\end{tikzpicture}
\end{center}

\begin{tcolorbox}[docTech, title={Document technique -- Figure 13 du sujet : caract\'eristiques en fonction du type de thermocouple}]
\footnotesize
\renewcommand{\arraystretch}{1.15}
\begin{center}
\begin{tabular}{@{}l p{1.65cm} p{1.75cm} c p{1.35cm} p{2.9cm}@{}}
\toprule
\textbf{Type} & \textbf{Plage de mesure conseill\'ee en \degree C} & \textbf{Alliages} & \textbf{Tol\'erance (\degree C)} & \textbf{Compa\-raison des prix} & \textbf{Fiabilit\'e, dur\'ee de vie, maintenance} \\
\midrule
T & $-200$ \`a $350$  & Cuivre et Constantan & $+/-1$   & Tr\`es \'elev\'e & R\'ep\'etabilit\'e exceptionnelle de $+/-0{,}1$ \degree C de $-200$ \`a $200$\degree C \\
K & $-10$ \`a $1\,100$ & Chromel et Alumel   & $+/-2{,}5$ & Faible          & Stabilit\'e moins satisfaisante que d'autres thermocouples \\
N & $0$ \`a $1250$     & Nicrosil et Nisil   & $+/-2{,}5$ & Tr\`es \'elev\'e & Stabilit\'e sup\'erieure aux autres thermocouples courants \\
J & $-30$ \`a $700$    & Fer et Constantan   & $+/-1{,}5$ & \'Elev\'e        & S'oxyde rapidement en milieu humide \\
E & $0$ \`a $800$      & Chromel et Constantan & $+/-2$ & \'Elev\'e        & Thermocouple \`a la FEM la plus \'elev\'ee \\
\bottomrule
\end{tabular}
\end{center}
\end{tcolorbox}

Le cahier des charges impose une \'etendue de mesure de \textbf{$-20$\degree C \`a $+200$\degree C}. L'armoire \'electrique est un \textbf{milieu sec} (local technique ventil\'e), et la temp\'erature de d\'eclenchement du ventilateur est de l'ordre de \SI{40}{\degreeCelsius}. Le co\^ut est un crit\`ere de s\'election explicite du dossier.

\begin{enumerate}[leftmargin=1.5em]
  \item \textbf{Relever} dans le document technique, pour chacun des cinq types de thermocouple, la plage de mesure conseill\'ee, et \textbf{indiquer} lesquels couvrent int\'egralement l'\'etendue $[-20\,;\,200]$\degree C impos\'ee par le cahier des charges.
  \item Parmi les types ainsi retenus, \textbf{justifier} le choix d'un thermocouple de \textbf{type J} en confrontant successivement les crit\`eres de tol\'erance, de prix et de tenue au milieu.
  \item Le thermocouple est un capteur \textbf{actif}. Rappeler ce que signifie ce qualificatif, et en d\'eduire pourquoi la cha\^ine d'acquisition ne peut pas se limiter \`a relier directement le capteur \`a l'entr\'ee analogique de l'automate.
\end{enumerate}

}{
\textbf{Corrig\'e.}

\begin{enumerate}[leftmargin=1.5em]
  \item \textbf{\'Etape 1 -- Extraction des donn\'ees.} Lecture directe de la colonne « plage de mesure conseill\'ee » et confrontation \`a l'exigence $[-20\,;\,200]$\degree C :
  \begin{center}
  \begin{tabular}{@{}lcc@{}}
  \toprule
  Type & Plage conseill\'ee (\degree C) & Couvre $[-20\,;\,200]$ ? \\
  \midrule
  T & $-200$ \`a $350$   & \textbf{oui} \\
  K & $-10$ \`a $1\,100$ & \textbf{non} ($-10 > -20$) \\
  N & $0$ \`a $1250$     & \textbf{non} ($0 > -20$) \\
  J & $-30$ \`a $700$    & \textbf{oui} \\
  E & $0$ \`a $800$      & \textbf{non} ($0 > -20$) \\
  \bottomrule
  \end{tabular}
  \end{center}
  C'est la \textbf{borne basse} qui \'elimine : pour mesurer jusqu'\`a $-20$\degree C, seuls les types \textbf{T} et \textbf{J} conviennent. Les types K, N et E ont une borne haute tr\`es confortable mais ne descendent pas assez bas.
  \item Il reste \`a d\'epartager le type T et le type J :
  \begin{itemize}
    \item \textbf{Tol\'erance} : le type T est nettement meilleur ($\pm 1$\degree C, avec une r\'ep\'etabilit\'e annonc\'ee de $\pm 0{,}1$\degree C entre $-200$ et $200$\degree C) que le type J ($\pm 1{,}5$\degree C). Mais l'application est un simple \textbf{seuil de ventilation} autour de \SI{40}{\degreeCelsius} : une incertitude de $\pm 1{,}5$\degree C sur l'instant d'enclenchement d'un ventilateur est sans cons\'equence. \textbf{L'exigence n'impose pas la pr\'ecision du type T.}
    \item \textbf{Prix} : le type T est class\'e « tr\`es \'elev\'e », le type J « \'elev\'e ». \`A performance suffisante des deux c\^ot\'es, le crit\`ere \'economique tranche en faveur du \textbf{type J}.
    \item \textbf{Tenue au milieu} : le seul d\'efaut r\'edhibitoire du type J est qu'il « s'oxyde rapidement en milieu humide ». L'armoire \'electrique \'etant un \textbf{milieu sec}, cette limitation ne s'applique pas ici.
  \end{itemize}
  \textbf{Le type J est donc retenu} : il couvre l'\'etendue de mesure, sa tol\'erance est suffisante au regard de l'usage, il est moins cher que le type T, et son unique faiblesse (l'humidit\'e) est hors du domaine d'emploi.

  \smallskip
  \emph{Remarque de m\'ethode.} Le raisonnement attendu est un \textbf{tri s\'equentiel} : d'abord \'eliminer sur le crit\`ere \emph{imp\'eratif} (l'\'etendue de mesure), puis arbitrer entre les survivants sur les crit\`eres \emph{pond\'er\'es} (pr\'ecision, co\^ut, environnement). Annoncer « type J » sans avoir \'elimin\'e K, N et E sur la borne $-20$\degree C ne r\'epond pas \`a la question pos\'ee.
  \item Un capteur \textbf{actif} \textbf{g\'en\`ere lui-m\^eme} la grandeur \'electrique de sortie \`a partir de l'\'energie pr\'elev\'ee sur le mesurande : il se comporte en \textbf{g\'en\'erateur} (ici une f.\,\'e.\,m. $e_{AB}$ par effet Seebeck) et ne n\'ecessite aucune source d'alimentation ext\'erieure. Il s'oppose au capteur \textbf{passif} (PT100, jauge), simple \textbf{imp\'edance variable} exigeant un circuit d'alimentation pour produire un signal.

  La f.\,\'e.\,m. d'un thermocouple est de l'ordre de \textbf{quelques millivolts} sur toute l'\'etendue de mesure. Elle est donc :
  \begin{itemize}
    \item \textbf{trop faible} devant la pleine \'echelle d'une entr\'ee analogique d'automate (typiquement \SIrange{0}{10}{\volt} ou \SIrange{4}{20}{\milli\ampere}) : num\'eris\'ee telle quelle, elle n'occuperait qu'une fraction infime du codage, d'o\`u une r\'esolution effective d\'erisoire ;
    \item \textbf{tr\`es vuln\'erable aux perturbations} sur la liaison capteur/armoire : quelques millivolts de mode commun induits par le c\^ablage de puissance voisin suffisent \`a noyer le signal utile.
  \end{itemize}
  Il faut donc interposer un \textbf{transmetteur} qui amplifie, d\'ecale puis convertit la tension en un \textbf{courant normalis\'e \SIrange{4}{20}{\milli\ampere}} (objet des exercices 3 et 4).

  \smallskip
  \textbf{Conclusion.} Le thermocouple de \textbf{type J} satisfait l'exigence d'\'etendue $[-20\,;\,200]$\degree C avec la tol\'erance requise pour une r\'egulation de ventilation, au co\^ut le plus bas compatible avec cette exigence, et sans exposition au seul milieu qui le d\'egraderait. Son caract\`ere actif et le tr\`es faible niveau de sa f.\,\'e.\,m. imposent en revanche un conditionnement complet en amont de l'automate : le choix du capteur \textbf{engage la cha\^ine d'acquisition enti\`ere}, il ne se r\'eduit pas \`a une lecture de catalogue.

  \smallskip
  \textit{Point de vigilance jury (rapport de jury CAPET 2024, analyse globale de l'\'epreuve \'ecrite disciplinaire) : \guillemotleft\ Le jury attend des r\'eponses d\'etaill\'ees et argument\'ees. Des r\'esultats donn\'es directement, sans calcul et sans justification ne r\'epondent pas aux attendus de l'\'epreuve. \guillemotright}
\end{enumerate}
}
\end{tcolorbox}

% =====================================================================
% Exercice 2 -- 2023, questions 46, 47 et 50
% =====================================================================
\begin{tcolorbox}[exoBox,
  title={Exercice 2 -- \difficulte{1} \hfill\normalfont\footnotesize 2023, \'epreuve \'ecrite disciplinaire appliqu\'ee, questions 46, 47 et 50}
]
\CAPETSHOW{
Les ch\^assis de surpression d'une serre horticole sont manoeuvr\'es par un motor\'educteur. Leur position angulaire $\alpha$ est mesur\'ee par un capteur \textbf{RPU} (\emph{Ridder Positioning Unit}), qui communique constamment la position au r\'egulateur climatique \`a l'aide d'un \textbf{signal \SIrange{4}{20}{\milli\ampere}}. Ce signal est g\'en\'er\'e par une sonde \`a effet Hall int\'egr\'ee au capteur, tournant sans contact devant un axe magn\'etique.

Le syst\`eme de fin de course RSU d\'efinit le \textbf{d\'ebattement angulaire m\'ecanique maximal} des ch\^assis : $-23\degree \leq \alpha \leq +23\degree$. Le capteur RPU est programm\'e \`a la mise en service en lui faisant apprendre un \textbf{trajet de mesurage} d\'elimit\'e par deux positions m\'emoris\'ees $SP_1$ et $SP_2$.

Dans un premier temps, le technicien programme le capteur avec $SP_1 = ES_{11}$ et $SP_2 = ES_{12}$, c'est-\`a-dire un trajet de mesurage confondu avec le d\'ebattement m\'ecanique maximal.

\begin{center}
\begin{tikzpicture}[>=Stealth, scale=0.9, every node/.style={scale=0.9}]
  \draw[->] (0,0) -- (7.2,0) node[right, font=\scriptsize] {$\alpha$ (\degree)};
  \draw[->] (0,0) -- (0,3.5) node[above, font=\scriptsize] {$I$ (mA)};
  \draw[thick, colExo] (0.6,0.6) -- (6.0,3.0);
  \fill[colExo] (0.6,0.6) circle (2pt);
  \fill[colExo] (6.0,3.0) circle (2pt);
  \draw[dotted] (0.6,0) -- (0.6,0.6) -- (0,0.6);
  \draw[dotted] (6.0,0) -- (6.0,3.0) -- (0,3.0);
  \draw[dotted] (3.3,0) -- (3.3,1.8) -- (0,1.8);
  \fill[colExo!60] (3.3,1.8) circle (1.6pt);
  \node[below, font=\scriptsize] at (0.6,-0.05) {$-23$};
  \node[below, font=\scriptsize] at (3.3,-0.05) {$0$};
  \node[below, font=\scriptsize] at (6.0,-0.05) {$+23$};
  \node[left, font=\scriptsize] at (-0.05,0.6) {$4$};
  \node[left, font=\scriptsize] at (-0.05,1.8) {?};
  \node[left, font=\scriptsize] at (-0.05,3.0) {$20$};
  \node[font=\scriptsize, align=center, colExo] at (4.95,1.05) {trajet de mesurage\\initial};
\end{tikzpicture}
\end{center}

\begin{enumerate}[leftmargin=1.5em]
  \item \textbf{Justifier} l'int\'er\^et de travailler avec la plage \SIrange{4}{20}{\milli\ampere} plut\^ot qu'avec la plage \SIrange{0}{20}{\milli\ampere}.
  \item \'Etablir l'expression de l'intensit\'e $I$ (en \si{\milli\ampere}) d\'elivr\'ee par le capteur en fonction de la position angulaire $\alpha$ (en degr\'es), puis calculer $I$ pour $\alpha = 0\degree$.
  \item En fonctionnement normal, la consigne d'ouverture des ch\^assis n'exc\`ede jamais \SI{50}{\percent}, ce qui correspond \`a $-23\degree \leq \alpha_{cons} \leq 0\degree$. Le technicien reprogramme alors $SP_2' = 0\degree$ (le trajet de mesurage devient plus court que le d\'ebattement m\'ecanique). \'Etablir la nouvelle expression de $I$ et \textbf{quantifier} le gain obtenu sur la r\'esolution de la mesure. Conclure.
\end{enumerate}

}{
\textbf{Corrig\'e.}

\begin{enumerate}[leftmargin=1.5em]
  \item Avec une plage \SIrange{4}{20}{\milli\ampere}, la valeur minimale du mesurande correspond \`a un courant \textbf{non nul} de \SI{4}{\milli\ampere}. C'est ce que l'on appelle le \textbf{z\'ero vivant} (\emph{live zero}).

  Cons\'equence directe : une \textbf{rupture de ligne} ou une \textbf{panne d'alimentation} du capteur donnent un courant voisin de \textbf{z\'ero}, valeur qui n'appartient pas \`a la plage de mesure utile. Le r\'egulateur peut donc \textbf{distinguer sans ambigu\"it\'e} une position basse l\'egitime ($\alpha = -23\degree \Rightarrow I = \SI{4}{\milli\ampere}$) d'une \textbf{d\'efaillance de la cha\^ine de mesure} ($I \approx 0$) : le z\'ero devient un indice de surveillance et permet de d\'eclencher une alarme.

  Avec une plage \SIrange{0}{20}{\milli\ampere}, $I = 0$ serait une valeur de mesure parfaitement valide : une rupture de c\^able serait interpr\'et\'ee comme « ch\^assis en but\'ee basse », erreur silencieuse et potentiellement dangereuse pour la culture.

  \emph{Avantage annexe} : la consommation minimale de \SI{4}{\milli\ampere} permet d'\textbf{alimenter l'\'electronique du capteur par la boucle elle-m\^eme} (transmetteur 2 fils), ce qu'une plage partant de \SI{0}{\milli\ampere} interdirait.
  \item La caract\'eristique est \textbf{lin\'eaire} et le trajet de mesurage couvre tout le d\'ebattement :
  \[ \alpha = -23\degree \;\longrightarrow\; I = \SI{4}{\milli\ampere} \qquad\text{et}\qquad \alpha = +23\degree \;\longrightarrow\; I = \SI{20}{\milli\ampere} \]
  La pente vaut :
  \[ k = \frac{20 - 4}{23 - (-23)} = \frac{16}{46} \approx 0{,}348~\mathrm{mA/\degree} \]
  et l'expression s'\'ecrit, en prenant le point $(-23\degree;\SI{4}{\milli\ampere})$ comme origine :
  \[ \boxed{\; I = 4 + \frac{16}{46}\,(\alpha + 23) \;} \qquad (I \text{ en mA},\; \alpha \text{ en degr\'es}) \]
  Pour $\alpha = 0\degree$ (ch\^assis \`a l'horizontale, milieu du d\'ebattement) :
  \[ I = 4 + \frac{16}{46}\times 23 = 4 + 8 = \boldsymbol{\SI{12}{\milli\ampere}} \]
  R\'esultat coh\'erent : le milieu de la plage angulaire donne bien le milieu de la plage de courant.
  \item Le trajet de mesurage est maintenant $[-23\degree\,;\,0\degree]$, et c'est sur cette \textbf{seule} plage que le capteur \'etale ses \SIrange{4}{20}{\milli\ampere} :
  \[ k' = \frac{20 - 4}{0 - (-23)} = \frac{16}{23} \approx 0{,}696~\mathrm{mA/\degree} \qquad\Longrightarrow\qquad \boxed{\; I = 4 + \frac{16}{23}\,(\alpha + 23) \;} \]
  La \textbf{sensibilit\'e} du capteur est donc multipli\'ee par :
  \[ \frac{k'}{k} = \frac{16/23}{16/46} = \frac{46}{23} = \boldsymbol{2} \]
  Autrement dit, un m\^eme \'ecart angulaire produit un \'ecart de courant \textbf{deux fois plus grand}. Pour une r\'esolution donn\'ee de la cha\^ine de num\'erisation en aval (quantum du convertisseur du r\'egulateur, exprim\'e en \si{\milli\ampere}), l'incertitude de positionnement angulaire est donc \textbf{divis\'ee par deux}. La zone $0\degree < \alpha \leq +23\degree$ sort en revanche du trajet de mesurage : elle n'est plus mesur\'ee (le capteur y sature \`a \SI{20}{\milli\ampere}).

  \smallskip
  \textbf{Conclusion.} Restreindre le trajet de mesurage \`a la seule plage r\'eellement exploit\'ee en fonctionnement normal double la r\'esolution effective de la mesure \textbf{sans changer de capteur ni toucher au r\'egulateur} : c'est un r\'eglage, pas un investissement. C'est la d\'emarche \`a engager d\`es lors qu'une exigence de pr\'ecision de positionnement n'est pas satisfaite avec le r\'eglage par d\'efaut. Le prix \`a payer -- la perte de mesure au-del\`a de $0\degree$ -- est ici acceptable puisque le cahier des charges garantit que la consigne n'exc\`ede jamais \SI{50}{\percent} d'ouverture ; il ne le serait pas si les ch\^assis devaient \^etre pilot\'es sur tout leur d\'ebattement.
\end{enumerate}
}
\end{tcolorbox}

% =====================================================================
% Exercice 3 -- 2024, questions E1, E2 et E3
% =====================================================================
\begin{tcolorbox}[exoBox,
  title={Exercice 3 -- \difficulte{2} \hfill\normalfont\footnotesize 2024, \'epreuve \'ecrite disciplinaire (option IE), questions E1, E2 et E3}
]
\CAPETSHOW{
On \'etudie l'\'etage d'\textbf{amplification} du transmetteur de temp\'erature qui met en forme la f.\,\'e.\,m. $e_{AB}$ du thermocouple de l'exercice 1. Cet \'etage est constitu\'e de \textbf{trois sous-\'etages en cascade} :
\[ e_{AB} \;\xrightarrow{\;\text{gain } G_1\;} \; V_1 \;\xrightarrow{\;G_2 = -100\;} \; V_2 \;\xrightarrow{\;G_3 = -1\;} \; V_3 \]

Le premier sous-\'etage est le montage ci-dessous, \`a deux amplificateurs op\'erationnels suppos\'es \textbf{parfaits} et fonctionnant en \textbf{r\'egime lin\'eaire}. Le thermocouple est connect\'e entre les deux entr\'ees non inverseuses ; $\varepsilon_1$ et $\varepsilon_2$ d\'esignent les tensions diff\'erentielles d'entr\'ee des deux AOP.

\begin{center}
\begin{circuitikz}[european, scale=0.92, every node/.style={scale=0.92}]
  % --- AOP : sorties en haut (y=4,4) et en bas (y=0) ; reseau resistif a droite
  \node [op amp, scale=0.85] (A1) at (3.0,4.4) {};
  \node [op amp, scale=0.85, yscale=-1] (A2) at (3.0,0.0) {};
  \node[font=\scriptsize, above] at (3.0,5.1) {AOP 1};
  \node[font=\scriptsize, below] at (3.0,-0.7) {AOP 2};
  % --- source eAB : segment vertical isole, relie aux deux entrees +
  \draw (0.5,4.4) to[battery1, l=$e_{AB}$, invert] (0.5,0.0);
  \draw (0.5,4.4) -- (A1.+);
  \draw (0.5,0.0) -- (A2.+);
  % tensions differentielles d'entree, decalees pour ne pas recouvrir les broches
  \draw[<->, thin] (2.08,4.68) -- (2.08,4.12);
  \node[left, font=\scriptsize] at (2.02,4.4) {$\varepsilon_1$};
  \draw[<->, thin] (2.08,0.28) -- (2.08,-0.28);
  \node[left, font=\scriptsize] at (2.02,0.0) {$\varepsilon_2$};
  % --- reseau R1 - R2 - R1 entre les deux sorties (bien espace)
  \draw (A1.out) -- (5.6,4.4);
  \draw (5.6,4.4) to[R, l_=$R_1$] (5.6,3.0);
  \draw (5.6,3.0) to[R, l_=$R_2$] (5.6,1.4);
  \draw (5.6,1.4) to[R, l_=$R_1$] (5.6,0.0);
  \draw (5.6,0.0) -- (A2.out);
  \node[circ] at (5.6,3.0) {};
  \node[circ] at (5.6,1.4) {};
  % --- contre-reaction : sorties des diviseurs vers les entrees inverseuses
  \draw (A1.-) -- (1.55,3.75) -- (1.55,3.0) -- (5.6,3.0);
  \draw (A2.-) -- (1.55,0.65) -- (1.55,1.4) -- (5.6,1.4);
  % --- VR2, mesuree aux bornes de R2
  \draw (5.6,3.0) -- (7.1,3.0);
  \draw (5.6,1.4) -- (7.1,1.4);
  \draw[<->, thin] (7.1,3.0) -- (7.1,1.4);
  \node[right, font=\scriptsize] at (7.2,2.2) {$V_{R2}$};
  % --- V1, mesuree entre les deux sorties d'AOP
  \draw (5.6,4.4) -- (8.5,4.4) node[circ]{};
  \draw (5.6,0.0) -- (8.5,0.0) node[circ]{};
  \draw[<->, thin] (8.85,4.4) -- (8.85,0.0);
  \node[right, font=\scriptsize] at (8.95,2.2) {$V_1$};
  \node[font=\scriptsize] at (4.6,-1.4) {$R_1 = \SI{1}{\kilo\ohm}$ \qquad $R_2 = \SI{834,4}{\ohm}$};
\end{circuitikz}
\end{center}

On donne la relation de sortie de ce premier sous-\'etage :
\[ V_1 = \frac{V_{R2}\,(2R_1 + R_2)}{R_2} \]

\begin{enumerate}[leftmargin=1.5em]
  \item Justifier que les deux AOP fonctionnent bien en r\'egime lin\'eaire, puis \textbf{\'ecrire} la loi des mailles impliquant $e_{AB}$, $\varepsilon_1$, $\varepsilon_2$ et $V_{R2}$. En d\'eduire que $e_{AB} = V_{R2}$.
  \item \textbf{Calculer} num\'eriquement le gain $G_1$ permettant d'\'ecrire $V_1 = G_1\,V_{R2} = G_1\,e_{AB}$.
  \item \textbf{Exprimer} la relation entre $e_{AB}$, $G_1$, $G_2$, $G_3$ et $V_3$, puis \textbf{exprimer} num\'eriquement $V_3$ en fonction de $e_{AB}$. Conclure sur l'ordre de grandeur obtenu et sur le r\^ole de la structure diff\'erentielle du premier \'etage.
\end{enumerate}

}{
\textbf{Corrig\'e.}

\begin{enumerate}[leftmargin=1.5em]
  \item \textbf{R\'egime lin\'eaire.} La sortie de chaque AOP est reboucl\'ee sur sa borne inverseuse ($-$) \`a travers le r\'eseau r\'esistif : la contre-r\'eaction est \textbf{n\'egative}, les AOP fonctionnent donc en \textbf{r\'egime lin\'eaire}. Les AOP \'etant suppos\'es parfaits, il vient :
  \[ \varepsilon_1 = \varepsilon_2 = \SI{0}{\volt} \qquad\text{et}\qquad i^+ = i^- = \SI{0}{\ampere} \]
  \textbf{Loi des mailles.} En parcourant la maille qui part de la borne $+$ de l'AOP 1, traverse le thermocouple, remonte par la branche $R_2$ puis se referme par l'entr\'ee de l'AOP 2 :
  \[ e_{AB} = \varepsilon_1 + V_{R2} - \varepsilon_2 \]
  Comme $\varepsilon_1 = \varepsilon_2 = 0$ :
  \[ \boxed{\; e_{AB} = V_{R2} \;} \]
  \emph{Interpr\'etation physique} : la f.\,\'e.\,m. du thermocouple se retrouve \textbf{int\'egralement} aux bornes de $R_2$. Le courant qui traverse la branche $R_1 - R_2 - R_1$ vaut donc $e_{AB}/R_2$ : c'est ce courant, \textbf{impos\'e par le capteur}, qui produit la tension de sortie. Les entr\'ees des AOP n'absorbant aucun courant ($i^+ = 0$), \textbf{le capteur n'est pas charg\'e} -- point essentiel pour une source de quelques millivolts \`a imp\'edance non n\'egligeable.
  \item En reportant $V_{R2} = e_{AB}$ dans la relation donn\'ee :
  \[ V_1 = \frac{V_{R2}\,(2R_1 + R_2)}{R_2} = \frac{2R_1 + R_2}{R_2}\;e_{AB} \]
  d'o\`u, num\'eriquement :
  \[ G_1 = \frac{2R_1 + R_2}{R_2} = \frac{2\times 1\,000 + 834{,}4}{834{,}4} = \frac{2\,834{,}4}{834{,}4} \approx \boldsymbol{3{,}397} \]
  \item Les trois sous-\'etages \'etant en cascade, les gains se \textbf{multiplient} :
  \[ V_3 = G_1\,G_2\,G_3\;e_{AB} \]
  \[ V_3 = 3{,}397 \times (-100) \times (-1) \times e_{AB} = \boldsymbol{340\;e_{AB}} \]
  Les deux inversions successives ($G_2 = -100$ puis $G_3 = -1$) se \textbf{compensent} : le gain global est \textbf{positif}, la tension $V_3$ varie donc dans le m\^eme sens que la temp\'erature. Un seul \'etage inverseur aurait donn\'e une caract\'eristique d\'ecroissante, incompatible avec la mise \`a l'\'echelle qui suit.

  \smallskip
  \textbf{Conclusion.} Le gain total de $340$ transforme une f.\,\'e.\,m. de quelques millivolts en une tension de quelques volts : pour $e_{AB} = \SI{10,8}{\milli\volt}$ (borne haute de la plage), $V_3 \approx 340 \times 10{,}8\times 10^{-3} \approx \SI{3,7}{\volt}$, valeur enfin compatible avec les \'etages suivants. Le choix d'un \textbf{premier \'etage \`a deux AOP en structure diff\'erentielle}, plut\^ot qu'un simple amplificateur non inverseur, est dict\'e par la nature du capteur : imp\'edance d'entr\'ee tr\`es \'elev\'ee (le thermocouple n'est pas charg\'e) et \textbf{r\'ejection du mode commun} (les perturbations induites de fa\c{c}on identique sur les deux fils de la liaison capteur/armoire sont \'elimin\'ees, seule la diff\'erence $e_{AB}$ est amplifi\'ee). Amplifier $340$ fois sans cette pr\'ecaution reviendrait \`a amplifier $340$ fois le bruit du c\^ablage de puissance voisin.
\end{enumerate}
}
\end{tcolorbox}

% =====================================================================
% Exercice 4 -- 2024, questions E5 et E10
% =====================================================================
\begin{tcolorbox}[exoBox,
  title={Exercice 4 -- \difficulte{2} \hfill\normalfont\footnotesize 2024, \'epreuve \'ecrite disciplinaire (option IE), questions E5 et E10}
]
\CAPETSHOW{
On poursuit l'\'etude du transmetteur de temp\'erature de l'exercice 3. Apr\`es l'\'etage d'amplification vient un \'etage de \textbf{d\'ecalage en tension} puis un \'etage de \textbf{conversion tension-courant} :
\[ e_{AB} \;\xrightarrow{\;\text{amplifier}\;}\; V_3 \;\xrightarrow{\;\text{d\'ecaler } (+V_e)\;}\; V_4 \;\xrightarrow{\;\text{convertir}\;}\; I_s \]

Le sujet donne, pour l'\'etage de d\'ecalage, la tension de r\'ef\'erence $V_e = \SI{1,34}{\volt}$ et la relation globale :
\[ V_4 = 337{,}7\,e_{AB} + 1{,}34 \qquad (V_4 \text{ en volts},\; e_{AB} \text{ en volts}) \]

L'\'etude de l'\'etage de conversion conduit \`a l'expression $I_s = V_4/R_5$, \textbf{ind\'ependante} de la r\'esistance de charge $R_{ana}$ (imp\'edance d'entr\'ee de la carte analogique de l'automate), avec $R_5 = \SI{250}{\ohm}$.

\begin{center}
\begin{tikzpicture}[>=Stealth, scale=0.9, every node/.style={scale=0.9}]
  \node[draw, minimum width=1.7cm, minimum height=0.9cm, align=center, font=\scriptsize] (a) at (1.5,0) {Amplifier\\$\times\,340$};
  \node[draw, minimum width=1.7cm, minimum height=0.9cm, align=center, font=\scriptsize] (b) at (4.2,0) {D\'ecaler\\$+\,V_e$};
  \node[draw, minimum width=1.9cm, minimum height=0.9cm, align=center, font=\scriptsize] (c) at (7.1,0) {Convertir\\$U \rightarrow I$};
  \draw[->] (-0.4,0) -- (a.west) node[midway, above, font=\scriptsize] {$e_{AB}$};
  \draw[->] (a.east) -- (b.west) node[midway, above, font=\scriptsize] {$V_3$};
  \draw[->] (b.east) -- (c.west) node[midway, above, font=\scriptsize] {$V_4$};
  \draw[->] (c.east) -- (9.1,0) node[midway, above, font=\scriptsize] {$I_s$};
  \node[font=\scriptsize] at (-0.55,-0.6) {mV};
  \node[font=\scriptsize] at (9.1,-0.6) {4-20 mA};
\end{tikzpicture}
\end{center}

Le capteur thermocouple de type J retenu \`a l'exercice 1 est \'etalonn\'e sur la plage $-20\degree$C $\rightarrow$ $200\degree$C ; le document technique DT10 du sujet donne les deux points caract\'eristiques extr\^emes :

\begin{center}
\begin{tabular}{@{}lcc@{}}
\toprule
Temp\'erature mesur\'ee (\degree C) & $-20$ & $200$ \\
Tension d\'elivr\'ee par le capteur (mV) & $e_{AB\,mini} = -1$ & $e_{AB\,maxi} = 10{,}8$ \\
\bottomrule
\end{tabular}
\end{center}

\begin{enumerate}[leftmargin=1.5em]
  \item \textbf{Calculer} $V_4$ pour les deux tensions extr\^emes de $e_{AB}$.
  \item En d\'eduire le r\^ole pr\'ecis du terme $+\,V_e = \SI{1,34}{\volt}$ : \`a quelle contrainte de la cha\^ine r\'epond-il, et que se passerait-il s'il \'etait omis ?
  \item \textbf{Calculer} $I_s$ pour les deux valeurs extr\^emes de $V_4$ et \textbf{conclure} quant \`a la conformit\'e du transmetteur de temp\'erature.
  \item L'expression $I_s = V_4/R_5$ ne fait pas intervenir $R_{ana}$. Expliquer en quoi cette propri\'et\'e est d\'eterminante pour une liaison industrielle.
\end{enumerate}

}{
\textbf{Corrig\'e.}

\begin{enumerate}[leftmargin=1.5em]
  \item Application directe de $V_4 = 337{,}7\,e_{AB} + 1{,}34$, avec $e_{AB}$ \textbf{exprim\'ee en volts} (attention \`a la conversion mV $\rightarrow$ V) :
  \[ e_{AB\,mini} = \SI{-1}{\milli\volt} \;\Longrightarrow\; V_4 = 337{,}7\times(-0{,}001) + 1{,}34 = -0{,}338 + 1{,}34 \approx \boldsymbol{\SI{1}{\volt}} \]
  \[ e_{AB\,maxi} = \SI{10,8}{\milli\volt} \;\Longrightarrow\; V_4 = 337{,}7\times 0{,}0108 + 1{,}34 = 3{,}647 + 1{,}34 \approx \boldsymbol{\SI{5}{\volt}} \]

  \emph{Remarque de rigueur.} Les valeurs exactes sont $\SI{1,0023}{\volt}$ et $\SI{4,987}{\volt}$ : le coefficient $337{,}7$ est lui-m\^eme un arrondi. Le coefficient qui donnerait \textbf{exactement} \SIrange{1}{5}{\volt} sur $[-1\,;\,10{,}8]$\,mV vaut $\frac{5-1}{(10{,}8+1)\times 10^{-3}} = \frac{4}{0{,}0118} \approx 339$, tr\`es proche du gain $340$ \'etabli \`a l'exercice 3. L'\'ecart de \SI{0,3}{\percent} sur la pleine \'echelle est sans cons\'equence pratique, mais il faut savoir d'o\`u il vient plut\^ot que de « retrouver » $1$ et $5$ par complaisance.
  \item L'\'etage d'amplification seul produit une tension $V_3$ qui, pour $e_{AB} = \SI{-1}{\milli\volt}$, vaut $340\times(-0{,}001) \approx \SI{-0,34}{\volt}$ : elle est \textbf{n\'egative}. Or l'\'etage suivant doit fournir un courant $I_s$ \textbf{strictement positif}, compris entre $4$ et \SI{20}{\milli\ampere}.

  Le terme $+\,V_e$ r\'ealise la \textbf{translation d'origine} (l'\emph{offset}) qui transforme la plage $[-0{,}34\,;\,3{,}65]$\,V, \`a cheval sur z\'ero, en la plage $[1\,;\,5]$\,V enti\`erement positive :
  \[ V_4 = \underbrace{340\,e_{AB}}_{\text{plage } [-0{,}34\,;\,+3{,}65]\,\mathrm{V}} + \underbrace{1{,}34}_{\text{d\'ecalage}} \;\longrightarrow\; \text{plage } [1\,;\,5]\,\mathrm{V} \]
  Sans ce d\'ecalage, toute la partie n\'egative de l'\'etendue de mesure (les temp\'eratures les plus basses) serait \textbf{perdue} : le convertisseur tension-courant ne pouvant d\'elivrer de courant n\'egatif, il saturerait \`a \SI{0}{\milli\ampere} et toutes ces temp\'eratures seraient confondues. Le d\'ecalage est donc la traduction \'electronique du \textbf{z\'ero vivant} de la boucle \SIrange{4}{20}{\milli\ampere} (cf. exercice 2).
  \item Avec $I_s = V_4/R_5$ et $R_5 = \SI{250}{\ohm}$ :
  \[ V_4 = \SI{1}{\volt} \;\Longrightarrow\; I_s = \frac{1}{250} = \SI{4e-3}{\ampere} = \boldsymbol{\SI{4}{\milli\ampere}} \]
  \[ V_4 = \SI{5}{\volt} \;\Longrightarrow\; I_s = \frac{5}{250} = \SI{20e-3}{\ampere} = \boldsymbol{\SI{20}{\milli\ampere}} \]

  \smallskip
  \textbf{Conclusion.} La cha\^ine compl\`ete transforme bien l'\'etendue de mesure $[-20\,;\,200]$\degree C du thermocouple de type J en la plage de sortie normalis\'ee $[4\,;\,20]$\,\si{\milli\ampere} attendue par le module d'entr\'ees analogiques de l'automate : \textbf{le transmetteur est conforme}. Les trois \'etages se lisent maintenant comme les trois op\'erations d'une \textbf{mise \`a l'\'echelle affine} : l'amplification fixe la \textbf{pente} ($\times 340$, du millivolt au volt), le d\'ecalage fixe l'\textbf{origine} ($+\SI{1,34}{\volt}$, pour rendre la plage positive), la conversion fixe la \textbf{nature du signal} transmis ($\div\,\SI{250}{\ohm}$, de la tension au courant). Aucun des trois n'est optionnel, et chacun se dimensionne \`a partir des deux points extr\^emes du cahier des charges.
  \item La grandeur transmise sur la ligne est un \textbf{courant}, et ce courant est \textbf{impos\'e} par le montage ind\'ependamment de la charge : le convertisseur se comporte en \textbf{source de courant} (tant que l'AOP ne sature pas, c'est-\`a-dire dans sa plage de conformit\'e). Trois cons\'equences industrielles majeures :
  \begin{itemize}
    \item la \textbf{chute de tension dans les conducteurs} (r\'esistance de ligne, qui s'ajoute \`a $R_{ana}$ sur des dizaines ou des centaines de m\`etres) n'affecte pas la valeur transmise -- alors qu'elle fausserait imm\'ediatement une transmission en tension ;
    \item la valeur ne d\'epend pas de l'imp\'edance d'entr\'ee exacte du r\'ecepteur : on peut changer de carte d'entr\'ees automate sans r\'e\'etalonner le transmetteur ;
    \item la boucle de courant est \textbf{intrins\`equement immunis\'ee} contre les perturbations \'electromagn\'etiques induites en mode commun sur la paire, qui n'injectent pas de courant dans une boucle ferm\'ee \`a faible imp\'edance.
  \end{itemize}
  C'est pr\'ecis\'ement ce qui justifie que le standard \SIrange{4}{20}{\milli\ampere} se soit impos\'e en instrumentation industrielle, l\`a o\`u le signal doit parcourir une distance notable entre le capteur et l'automate.
\end{enumerate}
}
\end{tcolorbox}

% =====================================================================
% Exercice 5 -- INEDIT -- pont de Wheatstone PT100, validite de dR << R0
% =====================================================================
\begin{tcolorbox}[exoBox,
  title={Exercice 5 -- \difficulte{3} \hfill\normalfont\footnotesize exercice in\'edit}
]
\CAPETSHOW{
Une sonde \textbf{PT100} mesure la temp\'erature $T$ d'un fluide caloporteur sur l'\'etendue $[0\,;\,150]$\degree C. C'est un capteur \textbf{passif} : sa r\'esistance suit la loi
\[ R_T = R_0\,(1 + \alpha T) \qquad\text{avec}\qquad R_0 = \SI{100}{\ohm} \;(\text{\`a } 0\degree\mathrm{C}), \quad \alpha = \SI{3,85e-3}{\per\degreeCelsius} \]

Elle est ins\'er\'ee dans un \textbf{pont de Wheatstone en quart de pont} (une seule branche active, les trois autres \'etant des r\'esistances fixes de valeur $R_0$), aliment\'e sous $V_{cc} = \SI{5}{\volt}$. La tension de d\'es\'equilibre $V_{out} = V_A - V_B$ attaque ensuite un amplificateur d'instrumentation de gain $K$, dont la sortie alimente un CAN d'entr\'ee \SIrange{0}{3}{\volt}.

\begin{center}
\begin{circuitikz}[european, scale=0.9, every node/.style={scale=0.9}]
  \draw (0,2.7) to[R, l_=$R_0$] (-1.9,1.35);
  \draw (-1.9,1.35) to[R, l_=$R_0$] (0,0);
  \draw (0,2.7) to[R, l=$R_0$] (1.9,1.35);
  \draw (1.9,1.35) to[R, l=$R_T$] (0,0);
  \node[circ] at (0,2.7) {};
  \node[circ] at (0,0) {};
  \node[circ] at (-1.9,1.35) {};
  \node[circ] at (1.9,1.35) {};
  \draw (0,2.7) -- (0,3.25) node[above, font=\scriptsize] {$+V_{cc} = \SI{5}{\volt}$};
  \draw (0,0) -- (0,-0.55) node[ground]{};
  \node[above left, font=\scriptsize] at (-1.95,1.4) {$B$};
  \node[above right, font=\scriptsize] at (1.95,1.4) {$A$};
  % tension de desequilibre, annotee au-dessus du centre du pont (hors des bras)
  \node[font=\scriptsize] at (0,1.35) {$V_{out} = V_A - V_B$};
  % liaisons vers ampli instrumentation
  \draw (1.9,1.35) -- (3.1,1.35) -- (3.1,1.75) -- (4.15,1.75);
  \draw (-1.9,1.35) -- (-2.9,1.35) -- (-2.9,-1.5) -- (3.55,-1.5) -- (3.55,0.95) -- (4.15,0.95);
  \node[draw, minimum width=1.9cm, minimum height=1.1cm, align=center, font=\scriptsize] (ai) at (5.15,1.35) {Ampli\\instrum.\\gain $K$};
  \node[font=\scriptsize] at (3.95,1.98) {$+$};
  \node[font=\scriptsize] at (3.95,0.7) {$-$};
  \draw[->] (ai.east) -- (7.1,1.35) node[right, font=\scriptsize, align=left] {vers CAN\\$[0\,;\,3]$ V};
\end{circuitikz}
\end{center}

\begin{enumerate}[leftmargin=1.5em]
  \item \'Etablir l'\textbf{expression exacte} de $V_{out} = V_A - V_B$ en fonction de $V_{cc}$, $R_0$ et $\Delta R = R_T - R_0$. Montrer que, \textbf{sous l'hypoth\`ese $\Delta R \ll R_0$}, elle se simplifie en $V_{out} \approx \dfrac{V_{cc}}{4}\cdot\dfrac{\Delta R}{R_0}$.
  \item Calculer $\Delta R$ et le rapport $\Delta R/R_0$ pour $T = \SI{10}{\degreeCelsius}$ puis pour $T = \SI{150}{\degreeCelsius}$. \textbf{L'hypoth\`ese de lin\'earisation est-elle l\'egitime sur toute l'\'etendue de mesure ?}
  \item Calculer, pour $T = \SI{150}{\degreeCelsius}$, la valeur de $V_{out}$ donn\'ee par la formule lin\'earis\'ee puis par l'expression exacte. Chiffrer l'\textbf{erreur relative} commise en utilisant la formule lin\'earis\'ee.
  \item D\'eterminer le gain $K$ de l'amplificateur d'instrumentation pour que la pleine \'echelle du CAN ($\SI{3}{\volt}$) soit atteinte \`a $T = \SI{150}{\degreeCelsius}$. Conclure sur la validit\'e de la cha\^ine ainsi con\c{c}ue.
\end{enumerate}

}{
\textbf{Corrig\'e.}

\begin{enumerate}[leftmargin=1.5em]
  \item Chaque branche du pont est un \textbf{diviseur de tension} entre $+V_{cc}$ et la masse. Les entr\'ees de l'amplificateur d'instrumentation n'absorbant aucun courant, les deux diviseurs sont ind\'ependants :
  \[ V_A = V_{cc}\,\frac{R_T}{R_0 + R_T} = V_{cc}\,\frac{R_0 + \Delta R}{2R_0 + \Delta R} \qquad\qquad V_B = V_{cc}\,\frac{R_0}{R_0 + R_0} = \frac{V_{cc}}{2} \]
  D'o\`u l'\textbf{expression exacte} :
  \[ V_{out} = V_A - V_B = V_{cc}\left(\frac{R_0 + \Delta R}{2R_0 + \Delta R} - \frac{1}{2}\right) = V_{cc}\,\frac{2(R_0 + \Delta R) - (2R_0 + \Delta R)}{2\,(2R_0 + \Delta R)} \]
  \[ \boxed{\; V_{out} = \frac{V_{cc}}{2}\cdot\frac{\Delta R}{2R_0 + \Delta R} \;} \]
  \textbf{Lin\'earisation.} Si et seulement si $\Delta R \ll R_0$, le terme $\Delta R$ est n\'egligeable devant $2R_0$ au d\'enominateur :
  \[ V_{out} \approx \frac{V_{cc}}{2}\cdot\frac{\Delta R}{2R_0} = \boxed{\;\frac{V_{cc}}{4}\cdot\frac{\Delta R}{R_0}\;} \]
  Il faut mesurer ce que l'on n\'eglige : l'approximation consiste \`a remplacer $2R_0 + \Delta R$ par $2R_0$, elle \textbf{surestime} donc $V_{out}$ d'un facteur $1 + \dfrac{\Delta R}{2R_0}$.
  \item $\Delta R = R_T - R_0 = R_0\,\alpha\,T$, soit num\'eriquement $\Delta R = 100 \times 3{,}85\times 10^{-3}\times T = 0{,}385\,T$ (en \si{\ohm}) :
  \begin{center}
  \begin{tabular}{@{}lccc@{}}
  \toprule
  $T$ & $\Delta R$ & $R_T = R_0 + \Delta R$ & $\Delta R / R_0$ \\
  \midrule
  \SI{10}{\degreeCelsius}  & \SI{3,85}{\ohm}  & \SI{103,85}{\ohm} & $0{,}0385$ \\
  \SI{150}{\degreeCelsius} & \SI{57,75}{\ohm} & \SI{157,75}{\ohm} & $\boldsymbol{0{,}5775}$ \\
  \bottomrule
  \end{tabular}
  \end{center}
  \`A \SI{10}{\degreeCelsius}, $\Delta R/R_0 \approx \SI{3,9}{\percent}$ : l'hypoth\`ese est raisonnable. Mais \`a \SI{150}{\degreeCelsius}, $\Delta R$ vaut \textbf{plus de la moiti\'e} de $R_0$ : \textbf{l'hypoth\`ese $\Delta R \ll R_0$ est manifestement fausse} en haut de l'\'etendue de mesure. C'est une diff\'erence de nature avec une \textbf{jauge de contrainte}, dont la variation relative reste de l'ordre de $10^{-4}$ et pour laquelle la lin\'earisation est toujours l\'egitime : \textbf{la formule $V_{cc}/4 \cdot \Delta R/R_0$ ne se transpose pas telle quelle d'un capteur \`a l'autre.}
  \item Pour $T = \SI{150}{\degreeCelsius}$, $\Delta R = \SI{57,75}{\ohm}$ :
  \[ V_{out}^{lin} = \frac{V_{cc}}{4}\cdot\frac{\Delta R}{R_0} = \frac{5}{4}\times\frac{57{,}75}{100} = 1{,}25\times 0{,}5775 \approx \boldsymbol{\SI{722}{\milli\volt}} \]
  \[ V_{out}^{exact} = \frac{V_{cc}}{2}\cdot\frac{\Delta R}{2R_0 + \Delta R} = \frac{5}{2}\times\frac{57{,}75}{200 + 57{,}75} = 2{,}5\times\frac{57{,}75}{257{,}75} \approx \boldsymbol{\SI{560}{\milli\volt}} \]
  Erreur relative de la formule lin\'earis\'ee :
  \[ \frac{V_{out}^{lin} - V_{out}^{exact}}{V_{out}^{exact}} = \frac{721{,}9 - 560{,}1}{560{,}1} \approx \boldsymbol{+\SI{28,9}{\percent}} \]
  On v\'erifie la coh\'erence avec l'analyse du 1. : le facteur de surestimation pr\'evu vaut $1 + \frac{\Delta R}{2R_0} = 1 + \frac{57{,}75}{200} = 1{,}289$, soit exactement les \SI{28,9}{\percent} trouv\'es. Une erreur de pr\`es de \SI{30}{\percent} sur la pleine \'echelle est \textbf{r\'edhibitoire} pour une mesure de temp\'erature.
  \item Le gain doit \^etre calcul\'e sur la valeur \textbf{exacte} de $V_{out}$ en haut d'\'echelle, sous peine de reproduire l'erreur pr\'ec\'edente :
  \[ K = \frac{V_{CAN\,max}}{V_{out}^{exact}(150\degree\mathrm{C})} = \frac{3}{0{,}5601} \approx \boldsymbol{5{,}36} \]
  \emph{Contre-exemple \`a ne pas commettre} : avec la formule lin\'earis\'ee on aurait trouv\'e $K = 3/0{,}7219 \approx 4{,}16$, gain \textbf{trop faible de \SI{22}{\percent}}. La sortie plafonnerait alors \`a $4{,}16\times 0{,}5601 \approx \SI{2,33}{\volt}$ \`a \SI{150}{\degreeCelsius} : plus de \SI{22}{\percent} de la dynamique du CAN resterait inutilis\'ee, d\'egradant d'autant la r\'esolution effective.

  \smallskip
  \textbf{Conclusion.} Avec $K \approx 5{,}4$, la cha\^ine couvre bien $[0\,;\,3]$\,V sur $[0\,;\,150]$\degree C \textbf{aux deux extr\'emit\'es}, mais elle reste \textbf{intrins\`equement non lin\'eaire} entre les deux : le pont de Wheatstone associ\'e \`a une PT100 sur une aussi large \'etendue ne fournit pas une tension proportionnelle \`a la temp\'erature. La conformit\'e de la cha\^ine ne peut donc pas \^etre prononc\'ee sur les seuls points extr\^emes -- il faut, soit \textbf{lin\'eariser par calcul} en aval (le traitement num\'erique inverse la relation exacte, solution retenue en pratique dans les automates et transmetteurs num\'eriques), soit \textbf{restreindre l'\'etendue de mesure} \`a un domaine o\`u $\Delta R \ll R_0$ est v\'erifi\'ee (par exemple $[0\,;\,20]$\degree C, o\`u l'erreur de lin\'earisation reste sous \SI{4}{\percent}). \'Ecrire $V_{out} = \frac{V_{cc}}{4}\frac{\Delta R}{R_0}$ sans jamais v\'erifier l'hypoth\`ese qui la fonde est l'erreur de fond \`a \'eviter.

  \smallskip
  \textit{Point de vigilance jury (rapport de jury CAPET 2025, commentaire de la partie 5) : \guillemotleft\ L'int\'er\^et d'un pont de Wheatstone et l'expression simplifi\'ee de la tension de sortie ont \'et\'e trait\'es par moins de vingt pour cent des candidats. \guillemotright}
\end{enumerate}
}
\end{tcolorbox}

% =====================================================================
% Exercice 6 -- 2022, question 24
% =====================================================================
\begin{tcolorbox}[exoBox,
  title={Exercice 6 -- \difficulte{3} \hfill\normalfont\footnotesize 2022, \'epreuve \'ecrite disciplinaire (option IE), question 24}
]
\CAPETSHOW{
Un module de t\'el\'etransmission TRT200, install\'e dans un poste de purge de r\'eseau vapeur, surveille l'\'etat de sa propre alimentation. Celle-ci est assur\'ee par un \textbf{module Peltier} (r\'ecup\'eration d'\'energie thermique sur la conduite) qui recharge des accumulateurs.

La tension $U_{Peltier}$ d\'elivr\'ee par le module est ramen\'ee \`a l'entr\'ee analogique \textbf{ADC2} du microcontr\^oleur PIC~18F8723 au travers d'un \textbf{pont diviseur} form\'e de deux r\'esistances \textbf{\'egales} $R_4$ et $R_5$ ; la tension appliqu\'ee sur ADC2 vaut donc la \textbf{moiti\'e} de la tension du module Peltier. Une diode $D_1$ prot\`ege l'entr\'ee contre les tensions n\'egatives.

Le convertisseur analogique-num\'erique interne au microcontr\^oleur est un \textbf{CAN 12 bits}, dont la tension de r\'ef\'erence vaut $V_{REF+} = \SI{3,5}{\volt}$ (et $V_{REF-} = \SI{0}{\volt}$).

\begin{center}
\begin{circuitikz}[european, scale=0.9, every node/.style={scale=0.9}]
  % --- pont diviseur R4 / R5, branche verticale bien etiree
  \draw (0.9,4.0) node[above, font=\scriptsize, align=center] {module Peltier\\$U_{Peltier}$} -- (0.9,3.6);
  \draw (0.9,3.6) to[R, l=$R_4$] (0.9,2.1);
  \draw (0.9,2.1) to[R, l=$R_5$] (0.9,0.6);
  \draw (0.9,0.6) -- (0.9,0.1) node[ground]{};
  \node[circ] at (0.9,2.1) {};
  % --- point milieu -> diode de protection (branche separee) -> CAN
  \draw (0.9,2.1) -- (2.9,2.1);
  \node[circ] at (2.9,2.1) {};
  \draw (2.9,2.1) to[D, l_=$D_1$] (2.9,0.6);
  \draw (2.9,0.6) -- (2.9,0.1) node[ground]{};
  \draw (2.9,2.1) -- (4.5,2.1) node[midway, above, font=\scriptsize] {ADC2};
  \node[draw, minimum width=2.6cm, minimum height=1.7cm, align=center, font=\scriptsize] (can) at (5.85,2.1) {CAN 12 bits\\PIC 18F8723\\$V_{REF+}{=}\SI{3,5}{\volt}$};
  \draw[->] (can.east) -- (7.85,2.1) node[right, font=\scriptsize, align=left] {$N$\\(0 \`a 4095)};
\end{circuitikz}
\end{center}

\begin{enumerate}[leftmargin=1.5em]
  \item \textbf{Calculer} le quantum $q$ du convertisseur analogique-num\'erique, en \si{\milli\volt} par LSB.
  \item En d\'eduire la valeur en sortie du CAN, exprim\'ee en base d\'ecimale, lorsque la tension aux bornes du module Peltier est de \SI{2}{\volt}.
  \item Le module TRT200 t\'el\'etransmet une alarme lorsque la sortie du CAN prend la valeur \texttt{0xAF} ($175$ en base d\'ecimale). Sachant que le cahier des charges impose que la tension aux bornes du module Peltier ne soit \textbf{pas inf\'erieure \`a \SI{0,9}{\volt}}, \textbf{justifier} le d\'eclenchement de cette alarme.
  \item Le pont diviseur divise la tension par deux \textbf{avant} la num\'erisation. Analyser l'effet de ce choix sur la r\'esolution de la mesure, et indiquer \`a quelle contrainte il r\'epond n\'eanmoins.
\end{enumerate}

}{
\textbf{Corrig\'e.}

\begin{enumerate}[leftmargin=1.5em]
  \item Le quantum est le pas de num\'erisation, c'est-\`a-dire la plus petite variation de tension d'entr\'ee qui fasse changer le code de sortie d'une unit\'e (LSB). Pour un convertisseur de $n$ bits, la plage $[V_{REF-}\,;\,V_{REF+}]$ est d\'ecoup\'ee en $2^n - 1$ intervalles :
  \[ q = \frac{V_{REF+} - V_{REF-}}{2^n - 1} = \frac{3{,}5 - 0}{2^{12} - 1} = \frac{3{,}5}{4\,095} \approx \SI{0,855e-3}{\volt} \approx \boldsymbol{\SI{0,85}{\milli\volt}\;/\;\mathrm{LSB}} \]
  \emph{Attention} : c'est bien $2^n - 1 = 4\,095$ et non $2^n = 4\,096$ qu'il faut \'ecrire au d\'enominateur, puisque le code varie de $0$ \`a $4\,095$, soit $4\,095$ intervalles pour $4\,096$ valeurs. L'\'ecart est faible (\SI{0,02}{\percent}) mais la confusion est syst\'ematiquement relev\'ee.
  \item Le pont diviseur \`a r\'esistances \'egales divise la tension par deux :
  \[ V_{ADC2} = \frac{U_{Peltier}}{2} = \frac{2}{2} = \SI{1}{\volt} \]
  Le code de sortie s'obtient en divisant la tension appliqu\'ee par le quantum, ou de fa\c{c}on \'equivalente par produit en croix sur la pleine \'echelle :
  \[ N = \frac{V_{ADC2}}{q} = V_{ADC2}\cdot\frac{4\,095}{V_{REF+}} = 1\times\frac{4\,095}{3{,}5} = \boldsymbol{1\,170} \]
  (valeur enti\`ere exacte ici, le CAN tronquant sinon \`a l'entier inf\'erieur).
  \item Il faut parcourir la cha\^ine \textbf{en sens inverse} : du code num\'erique vers la grandeur physique.

  \textbf{Du code vers la tension d'entr\'ee du CAN :}
  \[ V_{ADC2} = N\times q = 175\times\frac{3{,}5}{4\,095} \approx \SI{0,1496}{\volt} \approx \SI{0,15}{\volt} \]
  \textbf{De l'entr\'ee du CAN vers le module Peltier} (il faut « remonter » le pont diviseur, donc \textbf{multiplier} par deux) :
  \[ U_{Peltier} = 2\times V_{ADC2} = 2\times 0{,}15 = \boldsymbol{\SI{0,3}{\volt}} \]
  \textbf{Comparaison \`a l'exigence :} $\SI{0,3}{\volt} < \SI{0,9}{\volt}$. La tension d\'elivr\'ee par le module Peltier est donc \textbf{tr\`es en de\c{c}\`a} du minimum impos\'e : le d\'eclenchement de l'alarme t\'el\'etransmise est \textbf{pleinement justifi\'e}.

  \smallskip
  \emph{Pi\`ege de la cha\^ine inverse.} L'erreur classique est d'oublier le pont diviseur et de conclure sur les \SI{0,15}{\volt} de l'entr\'ee du CAN. Ici, elle ne change pas la conclusion (les deux valeurs sont sous le seuil), mais elle fausserait le diagnostic d\`es que la mesure s'approcherait de la limite : \textbf{le seuil du cahier des charges porte sur la grandeur physique, pas sur le signal conditionn\'e}. Chaque \'etage de la cha\^ine d'acquisition doit \^etre remont\'e explicitement.
  \item Le pont diviseur \textbf{divise par deux} la tension pr\'esent\'ee au CAN. \`A quantum constant (\SI{0,85}{\milli\volt} / LSB, fix\'e par $V_{REF+}$ et par les 12 bits), cela signifie qu'une variation donn\'ee de $U_{Peltier}$ ne produit plus que \textbf{moiti\'e moins} de codes : la r\'esolution rapport\'ee \`a la grandeur mesur\'ee passe de \SI{0,85}{\milli\volt} \`a \textbf{\SI{1,7}{\milli\volt} par LSB}. De plus, si $U_{Peltier}$ ne d\'epasse jamais quelques volts, une partie de la dynamique du CAN reste inutilis\'ee.

  Ce choix r\'epond n\'eanmoins \`a une \textbf{contrainte imp\'erative} : la tension d'entr\'ee d'un CAN ne doit \textbf{jamais} exc\'eder sa tension de r\'ef\'erence $V_{REF+} = \SI{3,5}{\volt}$, sous peine de saturation du code (\emph{clipping}, mesure fausse) voire de destruction de l'entr\'ee. Le module Peltier, selon son point de fonctionnement, peut d\'elivrer nettement plus de \SI{3,5}{\volt} ; le pont diviseur garantit que l'entr\'ee reste dans le domaine admissible sur toute la plage de fonctionnement. La diode $D_1$ compl\`ete ce r\^ole de protection du c\^ot\'e des tensions n\'egatives.

  \smallskip
  \textbf{Conclusion.} L'adaptation de niveau par pont diviseur est le \textbf{compromis classique du conditionnement} : on sacrifie d\'elib\'er\'ement une part de r\'esolution pour garantir l'\textbf{int\'egrit\'e} de l'entr\'ee de num\'erisation sur toute la plage de fonctionnement, protection qui prime sur la finesse de mesure d\`es lors que, comme ici, la fonction attendue est une \textbf{surveillance de seuil} ($U_{Peltier} \geq \SI{0,9}{\volt}$) et non une mesure de pr\'ecision. Dimensionner la r\'esolution au plus juste pour une simple fonction d'alarme serait l'erreur de conception inverse.

  \smallskip
  \emph{Remarque sur la source.} Le corrig\'e officiel de ce sujet compare la tension retrouv\'ee (\SI{0,3}{\volt}) \`a un seuil de \SI{0,4}{\volt} lu sur le diagramme des exigences, alors que l'\'enonc\'e du sujet \'enonce un minimum de \SI{0,9}{\volt}. Les deux seuils divergent dans les documents publi\'es ; la conclusion est identique dans les deux cas, puisque \SI{0,3}{\volt} est inf\'erieure \`a l'une comme \`a l'autre.

  \smallskip
  \textit{Point de vigilance jury (rapport de jury CAPET 2022, commentaires de l'\'epreuve \'ecrite disciplinaire) : \guillemotleft\ De nombreuses r\'eponses n'ont pas \'et\'e justifi\'ees ou reprennent sans aucune explication des donn\'ees issues de documentations techniques fournies. \guillemotright}
\end{enumerate}
}
\end{tcolorbox}

\CAPETENDFILE
